% !PRINTED 651
\weberpage{651}

Wenn also $\xi$ eine Zahl des Körpers $\Omega$ ist, die zu der
Gruppe $\mathfrak{A}$ gehört, so ist $\xi$ eine primitive Zahl des Körpers $b^{\mathrm{ten}}$
Grades $R(\xi)$, und in diesem Körper ist auch jede Zahl $\xi'$ enthalten,
die die Substitutionen der Gruppe $\mathfrak{A}'$ gestattet. Sind
also $\mathfrak{A}, \mathfrak{B}$ und $\mathfrak{A}', \mathfrak{B}'$ echte Theiler von $\mathfrak{G}$, so ist $R(\xi)$ von $\Omega$
verschieden, und $R(\xi)$ wird durch Adjunction von $\xi'$ nicht erweitert.
Es kann also auch nicht $\Omega$ aus $R(\xi)$ und $R(\xi')$ zusammengesetzt
werden, wodurch 3. bewiesen ist.

\begin{center}
\S.\,180.

\gesperrt{Die Resolventen.}
\end{center}

Nach dem, was jetzt bewiesen ist, können wir uns in den
weiteren Betrachtungen, die den Beweis des Satzes I. zum Ziele
haben, auf solche \gesperrt{Abel'sche} Körper beschränken, deren Grad
eine Primzahlpotenz und deren Gruppe cyklisch ist. Aus diesem
Grunde genügte es auch für die beabsichtigte Anwendung, dass
wir im vorigen Abschnitte uns auf die Betrachtung der Kreistheilungskörper
 $\Omega_m$ beschränkt haben, in denen $m$ eine Primzahlpotenz
 war. Wir behalten so viel als möglich die dort
gebrauchte Bezeichnung bei und setzen
\weq{1}{m = q^{\varkappa},}
worin $q$ eine Primzahl, $\varkappa$ ein positiver Exponent ist, und in dem
Falle, dass $q = 2$ ist, nehmen wir $\varkappa > 1$ an. Mit $n$ bezeichnen
wir jede durch $q$ nicht theilbare Zahl. Es ist $r$ eine primitive
$m^{\mathrm{te}}$ Einheitswurzel und $\Omega_m$ der Körper der rationalen Functionen
 von $r$.

Es sei nun $\Omega = R(x)$ ein einfacher \gesperrt{Abel'scher} Körper $m^{\mathrm{ten}}$
Grades, so dass $x$ Wurzel einer rationalen irreduciblen cyklischen
Gleichung $m^{\mathrm{ten}}$ Grades ist. Bezeichnen wir die Wurzeln dieser
Gleichung in geeigneter Reihenfolge mit
\weq{2}{x_0, x_1, x_2, \ldots, x_{m-1},}
und nehmen $x_m = x_0$, $x_{m+1} = x_1, \ldots$ an, so ist
\[
x_1 = \Phi(x_0), \quad x_2 = \Phi(x_1), \quad \ldots, \quad x_0 = \Phi(x_{m-1}),
\]
worin $\Phi$ eine ganze Function mit rationalen Coëfficienten bedeutet,
und die Gruppe des Körpers $\Omega$ besteht aus den Potenzen
der Substitution $(x_0, x_1)$ oder aus den cyklischen Permutationen
%CONT>>
